
\begin{corollaire}
\label{XIII.3.4}
Sean $k$ un cuerpo de característica $p\geqslant0$, $p'$ el conjunto de
los números primos distintos de~$p$ y $f\colon X\to k$ un morfismo
coherente. Supongamos que se cumple una de las condiciones siguientes:
\begin{enumerate}
\item[a)] $f$ es de tipo finito y los esquemas de tipo finito de dimensión
$\leqslant\dim X$ sobre una clausura algebraica de~$k$ son
desingularizables.
\item[b)] Los esquemas de tipo finito sobre una clausura algebraica de~$k$
son desingularizables.
\end{enumerate}
Entonces $f$ es universalmente localmente $1$-asférico para~$p'$.
\end{corollaire}

El caso a) resulta de~\Ref{XIII.3.3}. En el caso b), como la cuestión es
local sobre~$X$, puede suponerse que~$X$ es afín; por paso al límite
(SGA~4 XV~1.3), se reduce al caso en que~$X$ es de tipo finito sobre~$k$.

\begin{corollaire}
\label{XIII.3.5}
Sean
\marginpar{414}
$S$ un esquema irreducible con punto genérico~$s$ y $f\colon X\to S$ un
morfismo de presentación finita. Supongamos que los esquemas de tipo finito
de dimensión $\leqslant\dim X_s$ sobre una clausura
algebraica~$\overline{k}$ de~$\kres(s)$ son fuertemente desingularizables
\textup{(SGA~5 I~3.1.5)}. Si~$\LL$ designa el conjunto de los números primos
distintos de las características residuales de~$S$, puede encontrarse un
abierto no vacío~$S_1$ de~$S$ tal que, para toda especialización
$\bar{s}_1\to\bar{s}_2$ de puntos geométricos de~$S_1$, el morfismo de
especialización~\eqref{XIII.2.10}
$$
\pi_1^{\LL}(X_{\bar{s}_1})\to\pi_1^{\LL}(X_{\bar{s}_2})
$$
sea biyectivo.
\end{corollaire}

Según~\Ref{XIII.3.1} y~\Ref{XIII.3.3}, tras restringir~$S$ a un abierto no
vacío puede suponerse que~$f$ es localmente $1$-asférico para~$\LL$, y que,
para todo haz constante finito de~$\LL$-grupos~$F$ sobre~$X$, $(F,f)$ es
cohomológicamente propio en dimensión $\leqslant 1$. Resulta entonces
de~\Ref{XIII.1.14} que, para toda especialización
$\bar{s}_1\to\bar{s}_2$ de puntos geométricos de~$S_1$, el morfismo de
especialización
$$
\left(\R^1 f_* F\right)_{\bar{s}_2} \to \left(\R^1 f_*
F\right)_{\bar{s}_1}
$$
es biyectivo. El corolario no es sino la traducción de lo anterior en
términos de grupos fundamentales.

\section{Sucesiones exactas de homotopía}\label{XIII.4}\setcounter{subsection}{-1}
\subsection{}
\label{XIII.4.0}
Sean $X$ y $S$ dos esquemas conexos, $f\colon X\to S$ un morfismo, $a$ un
punto geométrico de~$X$ y $\LL$ un conjunto de números primos. Sea~$K$ el
núcleo del homomorfismo canónico $\pi_1(X,a)\to\pi_1(S,a)$ y~$N$ el menor
prosubgrupo normal de~$K$ tal que $K/N$ sea un grupo pro-$\LL$,
denotado por~$K^{\LL}$. Entonces $N$ es normal en~$\pi_1(X,a)$, y se denota por
$$
\pi_1'(X,a)
$$
\label{indnot:mj}\oldindexnot{$\pi_1'(X,a)$|hyperpage}%
el
\marginpar{415}
cociente de~$\pi_1(X,a)$ por~$N$. Si~$a$ es un punto geométrico de una
fibra geométrica~$X_{\bar{s}}$, los morfismos canónicos
$$
\pi_1(X_{\bar{s}},a)\to \pi_1(X,a) \to
\pi_1(S,a)
$$
permiten definir morfismos canónicos
$$
\pi_1^{\LL}(X_{\bar{s}},a)\lto{u} \pi_1'(X,a) \lto{v} \pi_1(S,a)
$$
Se tiene $vu=0$.
